\chapter{Design Documentation: Low Level}
\label{chap:design-low}

\section{Class Diagram}

The full class diagram below shows all major classes, inheritance relationships, aggregation, and key public interfaces.

\begin{figure}[H]
\centering
\includegraphics[width=\textwidth]{class_diagram.png}
\caption{Full class diagram for all sprints (see also Figure \ref{fig:class-diagram} in Process section)}
\label{fig:class-diagram-low}
\end{figure}

\subsection*{Player Class Hierarchy}

The central inheritance hierarchy separates game-logic sequencing from input method:

\begin{table}[H]
\centering
\footnotesize
\setlength{\tabcolsep}{3pt}
\begin{tabular}{|p{1.8cm}|p{2cm}|p{6cm}|}
\hline
\textbf{Class} & \textbf{Kind} & \textbf{Role} \\ \hline
\texttt{Player} & abstract Mono\-Be\-haviour & Owns all turn sequencing (\texttt{PlayTurn}), board-position delegation, knowledge tracking (\texttt{PlayerKnowledge}), and 147-cell notepad state. Declares abstract decision hook methods. \\ \hline
\texttt{HumanPlayer} & sealed : Player & Overrides every decision hook by calling \texttt{UiManager} and awaiting a callback via \texttt{WaitUntil}. \\ \hline
\texttt{AiPlayer} & sealed : Player & Overrides decision hooks: BFS-directed movement toward target room, knowledge-filtered suggestions, envelope-based accusation trigger. Comments describe further upgrade paths. \\ \hline
\end{tabular}
\end{table}

\subsection*{Key Public Decision Hook Signatures (\texttt{Player})}

\begin{table}[H]
\centering
\footnotesize
\setlength{\tabcolsep}{4pt}
\begin{tabular}{|p{5.2cm}|p{4.5cm}|}
\hline
\textbf{Method (abstract / public)} & \textbf{When called} \\ \hline
\texttt{ChooseTurnStart\-Action(\allowbreak canPassage, canSuggest, \allowbreak Action<TurnAction>)} & At turn start: roll / take passage / suggest-in-place \\ \hline
\texttt{ChooseMove(legal\-Moves, canEnterRoom, \allowbreak canUseSpecial, \allowbreak Action<Vector2Int,bool,\allowbreak bool>)} & Each movement step while walking \\ \hline
\texttt{ChooseDoor(doors, \allowbreak Action<Vector2Int>)} & When leaving a multi-door room \\ \hline
\texttt{ChooseSug\-gestion(roomID, \allowbreak chars, weapons, \allowbreak Action<Suggestion\-Input>)} & When the player must produce a suggestion \\ \hline
\texttt{ChooseAccusation\-FollowUp(\allowbreak Action<bool>)} & Immediately after suggestion resolves \\ \hline
\texttt{ChooseAccusation(\allowbreak roomID, Action<\allowbreak AccusationInput>)} & When \texttt{ChooseAccusationFollowUp} returned true \\ \hline
\texttt{ChooseCardToReveal(\allowbreak matches, Action<Card>)} & When refuting another player's suggestion \\ \hline
\texttt{ChooseRoomAction(\allowbreak canLeave, \allowbreak Action<RoomAction>)} & When in a room: stay and suggest in place, or leave through a door \\ \hline
\end{tabular}
\end{table}

\section{Key Class APIs}

\subsection*{\texttt{GameManagerScript} (Singleton)}

\begin{table}[H]
\centering
\footnotesize
\setlength{\tabcolsep}{4pt}
\begin{tabular}{|p{4.2cm}|p{5.5cm}|}
\hline
\textbf{Method} & \textbf{Description} \\ \hline
\texttt{CheckAccusation(\allowbreak character, weapon, \allowbreak room) : bool} & Checks three named cards against \texttt{Deck.\allowbreak GetEnvelope()} \\ \hline
\texttt{RunSuggestion(\allowbreak suggester, \allowbreak roomID) : \allowbreak IEnumerator} & Clockwise card-showing round; updates all players' \texttt{PlayerKnowledge}; shows revealed card to suggester via \texttt{UiManager.ShowRevealedCard} \\ \hline
\texttt{ResolveAccusation(\allowbreak accuser, \allowbreak Accusation\-Input) : \allowbreak IEnumerator} & Calls \texttt{CheckAccusation}; triggers \texttt{EndGame()} on win or sets \texttt{HasLost=true} on incorrect guess \\ \hline
\texttt{PlayersInTurnOrder\-StartingAfter(suggester) \allowbreak : IEnumerable\allowbreak <Player>} & Returns players clockwise from the suggester for card-showing \\ \hline
\texttt{GetNextFreeRoomSlot(\allowbreak roomID, excluding) \allowbreak : int} & Returns the next unoccupied visual slot in a room (1--6) \\ \hline
\texttt{MoveWeaponToRoom(\allowbreak weaponName, \allowbreak roomID)} & Repositions a weapon token to the specified room's spawn point \\ \hline
\texttt{DiceRoll() \allowbreak : GameObject} & Instantiates the physics dice and starts its roll coroutine \\ \hline
\texttt{EndGame()} & Sets \texttt{GameEnded=true}, destroys board/tokens, calls \texttt{GameState.GameOver()} \\ \hline
\end{tabular}
\end{table}

\subsection*{\texttt{BoardScript} (Singleton)}

\begin{table}[H]
\centering
\footnotesize
\setlength{\tabcolsep}{4pt}
\begin{tabular}{|p{4.2cm}|p{5.5cm}|}
\hline
\textbf{Member} & \textbf{Description} \\ \hline
\texttt{GetTile(x, y) : int} & Returns tile value; 0 for out-of-bounds \\ \hline
\texttt{SetTile(x, y, value)} & Writes a tile value; used by \texttt{TokenMover} to encode occupancy \\ \hline
\texttt{IsWalkable(x, y) \allowbreak : bool} & Returns true iff tile value is $\geq 1$ and $< 50$ (not wall, not player-start) \\ \hline
\texttt{FindPlayer(playerID) \allowbreak : Vector2Int} & Scans board for tile value == playerID; returns \texttt{(-1,-1)} if not found \\ \hline
\texttt{FindAllTiles(value) \allowbreak : List<Vector2Int>} & Returns all board positions with the given tile value (used for door sets) \\ \hline
\texttt{static RoomData : \allowbreak Dictionary<int,(string,\allowbreak Vector3)>} & Room ID $\to$ (name, world-space spawn position) \\ \hline
\texttt{const WALL=0, WALKABLE\-\_FLOOR=1, SPECIAL\-\_TILE=33, OCCUPIED=50} & Named tile-encoding constants \\ \hline
\end{tabular}
\end{table}

\subsection*{\texttt{TokenMover} (per-player MonoBehaviour)}

\begin{table}[H]
\centering
\footnotesize
\setlength{\tabcolsep}{4pt}
\begin{tabular}{|p{4.2cm}|p{5.5cm}|}
\hline
\textbf{Method} & \textbf{Description} \\ \hline
\texttt{InitialiseFromBoard(\allowbreak playerID)} & Locates the player's tile on \texttt{BoardScript}, sets grid position and world position \\ \hline
\texttt{GetLegalMoves() \allowbreak : List<Vector2Int>} & Returns orthogonal adjacent tiles that pass \texttt{BoardScript.IsWalkable} \\ \hline
\texttt{Step(Vector2Int delta)} & Moves one step: restores old tile value, encodes new tile with playerID, updates world position \\ \hline
\texttt{EnterRoom(\allowbreak roomID, slot)} & Moves token to the room's visual slot position; sets \texttt{CurrentRoomID} and \texttt{RoomSlot} \\ \hline
\texttt{LeaveThroughDoor(\allowbreak Vector2Int door)} & Exits room: places token on door tile, clears \texttt{CurrentRoomID}/\texttt{RoomSlot} \\ \hline
\end{tabular}
\end{table}

\subsection*{\texttt{PlayerKnowledge}}

\begin{table}[H]
\centering
\footnotesize
\setlength{\tabcolsep}{4pt}
\begin{tabular}{|p{4.2cm}|p{5.5cm}|}
\hline
\textbf{Method} & \textbf{Description} \\ \hline
\texttt{NoteOwnHand(\allowbreak owner, cards)} & Seeds the knowledge base with cards the player holds \\ \hline
\texttt{NoteShown(\allowbreak refuter, card)} & Records that \texttt{refuter} definitely holds \texttt{card} (suggester saw it) \\ \hline
\texttt{NoteCannotRefute(\allowbreak refuter, cards)} & Records that \texttt{refuter} holds none of \texttt{cards} (they passed) \\ \hline
\texttt{NoteRefutedSomething(\allowbreak refuter, cards)} & Records a disjunction: \texttt{refuter} holds at least one of \texttt{cards} (other observers) \\ \hline
\texttt{GetEnvelopeCandiates(\allowbreak allCards)} & Returns cards not definitively held by anyone---the set of possible envelope cards \\ \hline
\end{tabular}
\end{table}

\section{Object-Oriented Design Principles}

\subsection*{Single Responsibility}
Each class has one clearly bounded job:
\begin{itemize}
    \item \texttt{BoardScript} --- board state only (tile values, occupancy, walkability)
    \item \texttt{TokenMover} --- grid movement and room entry/exit only
    \item \texttt{Player} --- turn sequencing and knowledge tracking only; never touches the board directly
    \item \texttt{DiceRollScript} --- physics dice and face-detection only
    \item \texttt{UiManager} --- UI display and input collection only
    \item \texttt{Deck} / \texttt{Card} --- card data, dealing, and envelope management only
    \item \texttt{PlayerKnowledge} --- per-player card inference only
    \item \texttt{NotepadGrid} --- detective notebook toggle UI only
\end{itemize}

\subsection*{Inheritance and Polymorphism}
\texttt{Player} is an abstract MonoBehaviour that declares all turn-sequencing logic in concrete methods (\texttt{PlayTurn}, \texttt{DoMovementBudget}, \texttt{DoSuggestionThenAccusation}) and exposes abstract ``decision hook'' methods that the two subclasses override independently.
\begin{itemize}
    \item Template Method pattern: \texttt{Player.PlayTurn()} defines the skeleton of the turn; \texttt{HumanPlayer} and \texttt{AiPlayer} supply the input strategy by overriding each hook.
    \item Open/closed: adding a new player type (e.g.\ a network player) requires only a new subclass --- no changes to \texttt{GameManagerScript} or \texttt{Player.PlayTurn()}.
\end{itemize}

\subsection*{Singleton (Restricted to True Globals)}
\texttt{GameManagerScript}, \texttt{BoardScript}, \texttt{UiManager}, \texttt{GameState}, and \texttt{SceneController} use the Unity Singleton pattern because they represent unique global state any script may access. Per-player classes (\texttt{Player}, \texttt{TokenMover}, \texttt{DiceRollScript}) are instantiated per character.

\subsection*{Separation of Data and Presentation}
\texttt{Card} and \texttt{Deck} are pure C\# classes with no \texttt{MonoBehaviour} dependency. All visual representation is handled by \texttt{UiManager}. This keeps game logic testable independently of Unity's rendering pipeline.

\subsection*{Encapsulation}
Board state is private inside \texttt{BoardScript} --- all access is via \texttt{GetTile} / \texttt{SetTile} / \texttt{IsWalkable}. \texttt{TokenMover} writes back to the board via \texttt{SetTile}, so \texttt{BoardScript} is always the single source of truth for occupancy. \texttt{Deck}'s static dictionaries are private; access is only through \texttt{GetCharacterHand} and \texttt{GetEnvelope}.

\subsection*{Data-Driven Design}
All card names are loaded from \texttt{StreamingAssets/cards.json} at runtime via \texttt{JsonUtility}. Adding or renaming a character, weapon, or room requires only a JSON edit --- no code change, no recompile.

\subsection*{Callback-Driven UI Decoupling}
\texttt{UiManager} communicates results back to game logic via \texttt{System.Action} delegates; \texttt{Player} coroutines wait via \texttt{WaitUntil}. This means the turn coroutine has no knowledge of how the UI is built --- the same \texttt{Player.PlayTurn()} code drives both human and AI turns.

\section{Design vs.\ Implementation Deviations}

The table below records where the delivered codebase differs from the original design. Each deviation is documented transparently.

\begin{table}[H]
\centering
\footnotesize
\setlength{\tabcolsep}{4pt}
\begin{tabular}{|c|p{2.4cm}|p{2.8cm}|p{2.4cm}|p{2.8cm}|}
\hline
\textbf{\#} & \textbf{Designed} & \textbf{Delivered} & \textbf{Why} & \textbf{Where} \\ \hline
D1 & Standalone \texttt{Suggestion} & Inside \texttt{GameManager\-Script} & Late integration risk & \texttt{GameManager\-Script.cs} \\ \hline
D2 & Standalone \texttt{Accusation} & Inside \texttt{GameManager\-Script} & Same as D1 & \texttt{GameManager\-Script.cs} \\ \hline
D3 & Notebook-based AI deduction & BFS + knowledge-filtered & Core features in S4 & \texttt{AiPlayer.cs} \\ \hline
D4 & Full A* pathfinding & Legal-move navigation & Sprint 3 complexity & \texttt{TokenMover} \\ \hline
D5 & Single \texttt{TokenMovement\-Script} & Split into \texttt{Player} + subclasses & Separate concerns (S4) & 4 files \\ \hline
\end{tabular}
\end{table}

\section{Design Validation}

\subsection*{Sequence Diagrams}

The sequence diagrams are embedded in the sprint sections where the behaviour was first implemented:

\begin{itemize}
    \item \textbf{Player turn --- movement phase} (Figure \ref{fig:seq-movement}, Sprint 3): shows the interaction between \texttt{Player.PlayTurn()}, \texttt{TokenMover.GetLegalMoves()}, \texttt{HumanPlayer.ChooseMove()}, and \texttt{BoardScript}.
    \item \textbf{Suggestion resolution} (Figure \ref{fig:seq-suggestion}, Sprint 4): shows \texttt{GameManagerScript.RunSuggestion()} iterating players clockwise, calling \texttt{ChooseCardToReveal()} on the refuter, updating \texttt{PlayerKnowledge} on all observers, and calling \texttt{UiManager.ShowRevealedCard()} on the suggester.
\end{itemize}

\subsection*{Turn State Machine}

The turn state machine (Figure \ref{fig:state-turn}, Sprint 4) shows the states a player passes through within a single turn: \textit{Start} $\to$ \textit{ChooseAction} $\to$ \textit{Rolling/Moving/Passage} $\to$ \textit{InRoom} $\to$ \textit{Suggestion} $\to$ \textit{Accusation?} $\to$ \textit{End}.

% TODO: If the state machine diagram has been updated since the Sprint 4 image was generated
% (e.g. the ChooseTurnStartAction hook now also exposes a Suggest-in-place option),
% regenerate state_turn.puml and export a new state_turn.png to the overleaf/ folder.

\newpage
